\documentclass[11pt]{article}
\usepackage[margin=1in]{geometry}
\usepackage{amsmath,amssymb,amsthm}
\usepackage{booktabs}
\usepackage{hyperref}
\newtheorem{theorem}{Theorem}
\newtheorem{lemma}[theorem]{Lemma}
\newtheorem{proposition}[theorem]{Proposition}
\newtheorem{corollary}[theorem]{Corollary}
\theoremstyle{definition}
\newtheorem{definition}[theorem]{Definition}
\theoremstyle{remark}
\newtheorem{remark}[theorem]{Remark}
\newcommand{\F}{\mathbb{F}_2}
\newcommand{\wt}{\operatorname{wt}}
\newcommand{\dd}{\operatorname{d}}

\title{Disproof of Conjecture 00000007964:\\
binary linear quasi-perfect codes exist for every length $n\ge 2$}
\author{jilint777}
\date{2026-10-04}

\begin{document}
\sloppy
\maketitle

\begin{abstract}
Conjecture 00000007964 claims, among other things, that binary linear
quasi-perfect codes of length $n$ exist if and only if $n=2^k-1$ or $n$ lies
in a finite exceptional list. This is false. The conjecture defines
quasi-perfect codes as codes whose covering radius is one more than the packing
radius $e=\lfloor (d-1)/2\rfloor$ and which have a ``tight uniform shell
distribution''. We treat two readings: the standard one ($\rho=e+1$ only), and
the conjecture's stronger notion, where in addition any two words at the same
distance from the code have equally many codewords in every shell. Under
\emph{both} readings such codes exist for every length $n\ge 2$ and for no
length $n\le 1$. The witnesses are not degenerate:
for every $n\ge 3$ with $n+1$ not a power of two, the shortened Hamming code of
length $n$ is a quasi-perfect $[n,\,n-\lceil\log_2(n+1)\rceil,\,3]$ code. For
every $m\ge 2$, the extended Hamming code of length $2^m$ is quasi-perfect with
$d=4$. Both the extended Hamming codes and the once-shortened Hamming codes of
length $2^m-2$ ($d=3$) have uniform shell distributions. The set of lengths
outside $\{2^k-1\}$ is infinite, so no finite exceptional list can rescue the
clause. All three code families, the full length spectrum and the refutation
of the clause are formalized in Lean~4 (core library only). The Lean
refutation covers the uniform-shell notion (via even-weight codes), the
restrictions to $d\ge 3$ or $d\ge 4$, and the non-strict reading.
\end{abstract}

\section{The conjecture and the clause we refute}
The English statement reads:
\begin{quote}
\emph{Definition:} The classification extension of perfect codes: nontrivial
$q$-ary perfect codes (covering radius 1, balls fill the space) have known
types the Hamming codes and the binary/ternary Golay codes; define relaxed
quasi-perfect codes (Q-uniform ball measures) with tight uniform shell
distribution at covering radius $r+1$. \emph{Conjecture:} The gap constant of
the radius spectrum of quasi-perfect codes is at most 2 (covering efficiency
at least $1/2$); and binary linear quasi-perfect codes of length $n$ exist if
and only if $n$ is of type $2^k-1$ or in a finite exceptional list (the
Levenshtein set), completely enumerating the spectrum.
\end{quote}
The Chinese version says the same thing. The conjecture is a conjunction. Its
first conjunct (``gap constant of the radius spectrum'') is never defined, and
we do not discuss it. Refuting the second conjunct suffices:
\begin{equation}\label{eq:clause}
\begin{gathered}
\exists\ \text{finite } L\subset\mathbb{N}\ \ \forall n:\\
\bigl(\exists \text{ binary linear quasi-perfect code of length } n\bigr)
\iff \bigl(\exists k:\ n=2^k-1\bigr)\ \lor\ n\in L .
\end{gathered}
\end{equation}

\begin{definition}
A \emph{binary linear code} of length $n$ is an $\F$-subspace $C\subseteq\F^n$
with $|C|\ge 2$. Its \emph{minimum distance} is
$d=\min\{\dd(x,y): x\ne y\in C\}$ (Hamming distance), its \emph{packing
radius} is $e=\lfloor (d-1)/2\rfloor$, and its \emph{covering radius} is
$\rho=\max_{x\in\F^n}\min_{c\in C}\dd(x,c)$. The balls of radius $e$ around
codewords are disjoint and those of radius $\rho$ cover $\F^n$. Hence
$\rho\ge e$, and $C$ is \emph{perfect} if $\rho=e$. $C$ is
\emph{quasi-perfect} if $\rho=e+1$.
\end{definition}

This is the standard notion (for example, MacWilliams and Sloane). The
conjecture's own definition is \emph{stronger}. Its preamble asks for a
``tight uniform shell distribution at covering radius $r+1$'', where $r$ is
the packing radius. We formalize this as follows.

\begin{definition}\label{def:uniform}
For $x\in\F^n$ and $j\ge 0$ let $S_j(x)=\{c\in C:\dd(x,c)=j\}$ be the
$j$-th shell around $x$. The code $C$ is \emph{shell-uniform} if
$|S_j(x)|=|S_j(y)|$ for all $j$ whenever $\dd(x,C)=\dd(y,C)$. A
\emph{uniform-shell quasi-perfect} code is a quasi-perfect code that is
shell-uniform.
\end{definition}

This is the strongest reasonable reading: it constrains every shell, not only
the shells of radius $e$ and $e+1$ that enter the classical notion of a
uniformly packed code. We refute the clause under both the standard and the
uniform-shell reading (Section~\ref{sec:uniform} and Corollary~\ref{cor:uniform}).
Section~\ref{sec:readings} covers the remaining readings: the non-strict
version $\rho\le e+1$, a ``nontriviality'' demand $d\ge 3$ or $d\ge 4$,
restrictions on $k$, and the choice of the finite list.

\section{Syndromes}
We use the bitwise exclusive or $a\oplus b$ of nonnegative integers. It is
addition in the $\F$-vector space of finitely supported bit strings. In
particular it is associative and commutative, $a\oplus a=0$, $a\oplus 0=a$, and
$a\oplus b=0\iff a=b$.

Fix labels $\ell_1,\dots,\ell_n\in\mathbb{N}$, one per coordinate. The
\emph{syndrome} of $x\in\F^n$ is
\[
\sigma(x)=\bigoplus_{j:\,x_j=1}\ell_j .
\]
Then $\sigma(x+y)=\sigma(x)\oplus\sigma(y)$, so $\ker\sigma$ is a linear code.
This is just the code with parity-check matrix whose $j$-th column is the
binary expansion of $\ell_j$. Write $e_j$ for the $j$-th unit vector.

\begin{lemma}\label{lem:syn}
Suppose the labels are pairwise distinct.
\begin{enumerate}
\item If $\wt(x)=0$ then $\sigma(x)=0$. If $\wt(x)=1$, say $x=e_j$, then
$\sigma(x)=\ell_j$. If $\wt(x)=2$ then $\sigma(x)\neq 0$.
\item If all $\ell_j<2^m$ then $\sigma(x)<2^m$ for every $x$.
\item For $x\in\F^n$ and $y$ with $\sigma(y)=\sigma(x)$, the word $c=x+y$ lies
in $\ker\sigma$ and $\dd(x,c)=\wt(y)$.
\end{enumerate}
\end{lemma}
\begin{proof}
(1) $\sigma(e_i+e_j)=\ell_i\oplus\ell_j\ne 0$ for $\ell_i\ne\ell_j$.
(2) XOR does not create bits above position $m-1$.
(3) $\sigma(x+y)=\sigma(x)\oplus\sigma(x)=0$ and $x+c=y$.
\end{proof}

\begin{lemma}\label{lem:bit}
If $t=2^k\le s<2^{k+1}$ then $s\oplus t=s-t<t$.
\end{lemma}
\begin{proof}
Bit $k$ is the leading bit of $s$. XOR with $2^k$ clears it and leaves the
lower bits unchanged.
\end{proof}

\section{\texorpdfstring{Shortened Hamming codes: $d=3$, every length $n\ge 3$ with $n+1\ne 2^k$}{Shortened Hamming codes}}
For $n\ge 1$ let $\mathrm{SH}(n)=\ker\sigma\subseteq\F^n$ with labels
$\ell_j=j$ ($j=1,\dots,n$). For $n=2^m-1$ this is the Hamming code. For other
$n$ it is a shortened Hamming code.

\begin{theorem}\label{thm:sh}
Let $n\ge 3$ and suppose $n+1$ is not a power of $2$. Then $\mathrm{SH}(n)$ is
a binary linear code with minimum distance $d=3$ (so $e=1$) and covering
radius $\rho=2$. It is quasi-perfect.
\end{theorem}
\begin{proof}
\emph{Minimum distance.} Let $c\ne 0$ be a codeword. $\wt(c)\ne 1$, because
$\sigma(e_j)=j\ne 0$. $\wt(c)\ne 2$ by Lemma~\ref{lem:syn}(1). So
$\wt(c)\ge 3$. Since $1\oplus 2\oplus 3=0$, $e_1+e_2+e_3$ is a codeword of
weight $3$. For a linear code the minimum distance is the least nonzero weight,
so $d=3$.

\emph{Covering radius $\le 2$.} Let $k=\lfloor\log_2 n\rfloor$ and $t=2^k$, so
$t\le n<2t$. Let $x\in\F^n$ and $s=\sigma(x)$. By Lemma~\ref{lem:syn}(2) with
$m=k+1$ we have $s<2t$. We find $y$ with $\wt(y)\le 2$ and $\sigma(y)=s$:
\begin{itemize}
\item if $s=0$, take $y=0$;
\item if $1\le s\le n$, take $y=e_s$;
\item if $n<s<2t$, put $b=s\oplus t$. By Lemma~\ref{lem:bit}, $b<t\le n$.
Also $b\ne 0$, since otherwise $s=t\le n$. So $y=e_t+e_b$ is a word of weight
$2$ with $\sigma(y)=t\oplus b=s$.
\end{itemize}
By Lemma~\ref{lem:syn}(3), $c=x+y$ is a codeword with $\dd(x,c)\le 2$.

\emph{Covering radius $\ge 2$.} Since $n<2t$ and $n+1\ne 2t$, we have $n\le
2t-2$. Put $s^\ast=2t-1$, so $n<s^\ast<2t$. By the previous paragraph there is
$y^\ast$ with $\sigma(y^\ast)=s^\ast$. For every codeword $c$,
$\sigma(y^\ast+c)=s^\ast$. Since $s^\ast\ne 0$ and $s^\ast\notin\{1,\dots,n\}$,
Lemma~\ref{lem:syn}(1) gives $\wt(y^\ast+c)\notin\{0,1\}$. So
$\dd(y^\ast,c)\ge 2$ for all $c$, and $\rho=2=e+1$.
\end{proof}

For $n=2^m-1$ the same argument gives $\rho=1=e$: every $s<2^m$ is a label or
$0$. So $\mathrm{SH}(2^m-1)$ is the perfect Hamming code and not
quasi-perfect. The Lean file proves this as a sanity check
(\texttt{hamming\_not\_QP}). The dimension of $\mathrm{SH}(n)$ is
$n-\lceil\log_2(n+1)\rceil$, because the labels $1,2,4,\dots,2^k\le n$ make
$\sigma$ surjective onto $\{0,\dots,2^{k+1}-1\}$. This is not needed for the
disproof.

\section{\texorpdfstring{Extended Hamming codes: $d=4$, lengths $2^m$}{Extended Hamming codes}}
For $m\ge 2$ let $N=2^m$. Label the coordinates $0,1,\dots,N-1$ and let
\[
\mathrm{EH}(m)=\{x\in\F^N:\ \sigma(x)=0,\ \wt(x)\text{ even}\}.
\]
This is the extended Hamming code. It is linear because $\sigma$ and the
parity $\wt(x)\bmod 2$ are both linear.

\begin{theorem}\label{thm:eh}
$\mathrm{EH}(m)$ has $d=4$ (so $e=1$) and covering radius $2$, so it is
quasi-perfect. Moreover, every word at distance $2$ from the code has exactly
$N/2$ codewords at distance $2$. Every word at distance $1$ has exactly one
codeword at distance $1$ and none at distance $2$.
\end{theorem}
\begin{proof}
A nonzero codeword has even weight, and weight $2$ is excluded by
Lemma~\ref{lem:syn}(1). Since $0\oplus1\oplus2\oplus3=0$, $e_0+e_1+e_2+e_3$ is a
codeword of weight $4$. So $d=4$.

Let $x\in\F^N$ and $s=\sigma(x)<N$. If $\wt(x)$ is odd, $y=e_s$ has syndrome
$s$ and odd weight. If $\wt(x)$ is even and $s\ne 0$, $y=e_0+e_s$ has syndrome
$s$ and even weight. If $\wt(x)$ is even and $s=0$, $x$ is a codeword. In each
case $c=x+y$ is a codeword with $\dd(x,c)\le 2$. For $x=e_0+e_1$ and any
codeword $c$, $x+c$ has syndrome $1\neq 0$ and even weight, so
$\dd(x,c)\ge 2$. Hence $\rho=2$.

Shells: if $\dd(x,C)=2$ then $\wt(x)$ is even and $s=\sigma(x)\ne 0$. The
codewords at distance $2$ are the $x+e_a+e_b$ with $a\oplus b=s$, and there
are exactly $N/2$ unordered pairs $\{a,a\oplus s\}$. If $\dd(x,C)=1$, the
codeword at distance $1$ is unique, because $d=4>2$. By parity no codeword is
at even distance from $x$.
\end{proof}

The shell numbers depend only on the distance to the code. So
$\mathrm{EH}(m)$ is uniformly packed: its shell distribution at the covering
radius $e+1$ is as uniform as it can be.

\section{Even-weight and repetition codes}
\begin{proposition}\label{prop:ew}
For $n\ge 2$ the even-weight code $\mathrm{EW}(n)=\{x:\wt(x)\text{ even}\}$ has
$d=2$, $e=0$ and $\rho=1$, so it is quasi-perfect. Every word not in the code
has exactly $n$ codewords at distance $1$.
\end{proposition}
\begin{proof}
Nonzero codewords have weight $\ge 2$, and $e_1+e_2$ is a codeword. A word
of odd weight becomes a codeword by flipping any one of its $n$ bits, and it is
at odd distance from every codeword. So $\rho=1$, attained by $e_1$.
\end{proof}

For even $n=2m$ the repetition code $\{0,\mathbf 1\}$ has $d=2m$, $e=m-1$ and
$\rho=m$. Every word is within $\min(\wt x,n-\wt x)\le m$ of a codeword, and
weight-$m$ words are at distance $m$ from both codewords. So it is
quasi-perfect as well. This family is mentioned only for completeness and is
checked in \texttt{verify.py}.

\begin{proposition}\label{prop:short}
There is no binary linear quasi-perfect code of length $n\le 1$.
\end{proposition}
\begin{proof}
For $n=0$ there is only one word, so no code has two codewords. For $n=1$ the
only linear code with two codewords is $\F^1$. It has $\rho=0<e+1$.
\end{proof}

\begin{theorem}[length spectrum]\label{thm:spec}
A binary linear quasi-perfect code of length $n$ exists if and only if $n\ge 2$.
Codes with $d=3$ exist for every $n\ge 3$ with $n+1\ne 2^k$
(Theorem~\ref{thm:sh}). Codes with $d=4$ exist for every $n=2^m$, $m\ge 2$
(Theorem~\ref{thm:eh}).
\end{theorem}
\begin{proof}
Use Propositions~\ref{prop:ew} and \ref{prop:short}.
\end{proof}

\section{Uniform shell distributions}\label{sec:uniform}
For a linear code $C$, the map $c\mapsto x+c$ is a bijection from $S_j(x)$ onto
the set of weight-$j$ words of the coset $x+C$. Also $\dd(x,C)$ is the minimum
weight of $x+C$. Hence:

\begin{lemma}\label{lem:coset}
A linear code is shell-uniform if and only if any two cosets with the same
minimum weight have the same weight distribution. Suppose a coordinate
permutation $\pi$ maps $C$ onto itself. Then $\pi$ maps the coset $x+C$ onto
$\pi(x)+C$ and preserves weights. So if the group of such permutations acts
transitively on the cosets of each minimum weight, $C$ is shell-uniform.
\end{lemma}

\begin{theorem}\label{thm:uniform}
The following quasi-perfect codes are shell-uniform:
\begin{enumerate}
\item $\mathrm{EW}(n)$ for every $n\ge 2$;
\item $\mathrm{EH}(m)$ for every $m\ge 2$ ($d=4$, length $2^m$);
\item $\mathrm{SH}(2^m-2)$ for every $m\ge 3$ ($d=3$, length $2^m-2$).
\end{enumerate}
\end{theorem}
\begin{proof}
Identify a label $\ell<2^m$ with the vector of its binary digits in $\F^m$.
Then $\oplus$ is vector addition and $\sigma(x)=\sum_{j:x_j=1}\ell_j$.

(1) $\mathrm{EW}(n)$ has two cosets, itself (minimum weight $0$) and the
odd-weight words (minimum weight $1$). There is one coset of each minimum
weight. (In Lean the bijection $c\mapsto c+x+y$ is given explicitly.)

(2) The coordinates of $\mathrm{EH}(m)$ are labelled by all of $\F^m$. The map
$x\mapsto(\sigma(x),\wt(x)\bmod 2)$ is linear and surjective onto
$\F^m\times\F$, and its kernel is $\mathrm{EH}(m)$. So the cosets correspond to
the pairs $(s,p)$. By the proof of Theorem~\ref{thm:eh}, the minimum weights
are as follows: $(0,0)$ has minimum weight $0$; $(s,1)$ has minimum weight $1$
for every $s$; and $(s,0)$ with $s\neq 0$ has minimum weight $2$.

Every affine bijection $\varphi(v)=Av+b$ of $\F^m$, with $A\in GL(m,2)$,
permutes the coordinates. Let a word have support $T$. Its image has support
$\varphi(T)$, with syndrome $A\sigma+|T|\,b$ and the same parity. So $\varphi$
maps the coset $(s,p)$ to $(As+pb,\,p)$, and it preserves $\mathrm{EH}(m)$,
which is the coset $(0,0)$.
\begin{itemize}
\item The translations ($A=I$, $b=s+s'$) map $(s,1)$ to $(s',1)$.
\item For nonzero $s,s'$, a linear map $A$ with $As=s'$ maps $(s,0)$ to
$(s',0)$.
\end{itemize}
So the action is transitive on the cosets of each minimum weight.

(3) Let $n=2^m-2$ and $u=2^m-1$, the all-ones vector. The labels are
$\F^m\setminus\{0,u\}$. The unit vectors $1,2,\dots,2^{m-1}$ are labels, so
$\sigma$ is surjective and the cosets correspond to syndromes $s\in\F^m$. By
Theorem~\ref{thm:sh} and Lemma~\ref{lem:syn}, the minimum weights are: $s=0$
gives $0$; $s=u$ gives $2$; and every other $s$ (a label) gives $1$. Minimum
weights $0$ and $2$ each have a single coset.

Now let $s,s'\notin\{0,u\}$. Then $\{s,u\}$ and $\{s',u\}$ are both linearly
independent, so some $A\in GL(m,2)$ satisfies $Au=u$ and $As=s'$ (extend both
pairs to bases). $A$ fixes $0$ and $u$, so it permutes the labels. The induced
coordinate permutation multiplies syndromes by $A$, so it preserves the code
and maps coset $s$ to coset $s'$. Lemma~\ref{lem:coset} finishes the proof.
\end{proof}

This agrees with the classical literature on uniformly packed codes
(Semakov, Zinov'ev and Zaitsev; Goethals and van Tilborg), where extended and
once-shortened Hamming codes are standard examples. Our proof does not depend
on it. Shortened Hamming codes of other lengths are in general \emph{not}
shell-uniform. For example, in $\mathrm{SH}(5)$ the words $e_1$ and $e_2$ are
both at distance $1$ from the code, but only $e_1$ has a codeword at distance
$5$, namely $e_2+e_3+e_4+e_5$ (Lean: \texttt{SH5\_not\_shellUniform}).

\section{The disproof}
\begin{corollary}\label{cor:main}
Clause~\eqref{eq:clause} is false. In fact, already its ``only if'' direction
fails for every finite $L$. This holds even when only codes with $d\ge 3$ (or
$d\ge 4$) are counted, and even when perfect codes are also allowed (reading
$\rho\le e+1$).
\end{corollary}
\begin{proof}
Let $L$ be finite and $M=\max(L\cup\{0\})$. Put $n=2(M+2)$. Then $n\ge 4$,
$n\notin L$, and $n+1$ is odd and $\ge 5$, so it is not a power of $2$. By
Theorem~\ref{thm:sh}, $\mathrm{SH}(n)$ is quasi-perfect with $d=3$, but $n$ is
neither of the form $2^k-1$ nor in $L$. For the $d\ge 4$ version take
$n=2^{M+2}>M$. Then $n+1$ is odd and $>1$, and $\mathrm{EH}(M+2)$ is
quasi-perfect with $d=4$. A code with $\rho=e+1$ in particular satisfies
$\rho\le e+1$, so the non-strict reading is covered too.
\end{proof}

\begin{corollary}[uniform-shell reading]\label{cor:uniform}
Clause~\eqref{eq:clause} is also false when ``quasi-perfect'' means
uniform-shell quasi-perfect (Definition~\ref{def:uniform}). Such codes exist
exactly for $n\ge 2$, by $\mathrm{EW}(n)$ and Proposition~\ref{prop:short}.
Nontrivial ones exist for the infinitely many lengths $2^m-2$ ($d=3$) and $2^m$
($d=4$), and none of these lengths has the form $2^k-1$.
\end{corollary}
\begin{proof}
Use Theorem~\ref{thm:uniform}. Choose $n$ larger than $\max L$ exactly as in
Corollary~\ref{cor:main}: $n=2(M+2)$ for $\mathrm{EW}$, $n=2^{M+3}-2$ for
$\mathrm{SH}$, or $n=2^{M+2}$ for $\mathrm{EH}$.
\end{proof}

So the true spectrum is $\{n\ge 2\}$ under both readings, and its complement $\{0,1\}$ is
finite. The proposed ``spectrum'' $\{2^k-1\}\cup L$ has density zero. The
clause seems to confuse quasi-perfect codes with perfect codes. The lengths of
nontrivial binary perfect codes do lie in $\{2^k-1\}\cup\{23\}$, and the binary
Golay code supplies the exceptional length.

\section{Robustness of the interpretation}\label{sec:readings}
\begin{itemize}
\item \textbf{``Tight'' as ``nearly perfect''.} If ``tight'' is read in the
sense of Goethals and Snover (nearly perfect codes, i.e.\ codes meeting the
Johnson bound with equality), the conclusion is unchanged. For $e=1$ and $n$
even the Johnson bound reads $|C|\le 2^n/(n+2)$, and $\mathrm{SH}(2^m-2)$ has
$2^{n-m}=2^n/(n+2)$ codewords, so it meets the bound. These once-shortened
Hamming codes are the standard nearly perfect codes with $e=1$ (Goethals--Snover
1972; Lindstr\"om 1975). They give the infinitely many lengths $2^m-2$, none of the
form $2^k-1$. (On paper only.)
\item \textbf{Strict versus non-strict.} Our codes have $\rho=e+1$ exactly, so
they are quasi-perfect whether or not perfect codes are counted
(Lean: \texttt{onlyIfWeak\_false}).
\item \textbf{Trivial codes.} Someone could object to $d\le 2$ or to
repetition codes. The shortened Hamming codes ($d=3$, every admissible
$n\ge 4$) and the extended Hamming codes ($d=4$) are standard nontrivial codes
(Lean: \texttt{onlyIf3\_false}, \texttt{onlyIf4\_false},
\texttt{levenshteinAtLeast\_false} for every bound $D\le 4$ on $d$).
\item \textbf{Range of $k$, and the ``if'' direction.} We refute the ``only
if'' direction with $k$ arbitrary. This is the weakest form, and it implies
every variant that restricts $k$. The ``if'' direction is not used. (It holds
for $k\ge 2$ because of $\mathrm{EW}(2^k-1)$, and fails only at $n\in\{0,1\}$.)
\item \textbf{The exceptional list.} $L$ is arbitrary and finite, which
includes whatever the ``Levenshtein set'' is meant to be. The counterexample
lengths form an infinite set (all even $n\ge 4$ already suffice).
\item \textbf{The conjecture's uniform-shell notion.} This notion is handled
by Definition~\ref{def:uniform}, Theorem~\ref{thm:uniform} and
Corollary~\ref{cor:uniform}. In Lean, the even-weight family carries the
refutation (\texttt{EW\_shellUniform}, \texttt{UQP\_spectrum},
\texttt{conjecture\_00000007964\_false\_uniform}). The nontrivial uniform
families $\mathrm{SH}(2^m-2)$ ($d=3$) and $\mathrm{EH}(m)$ ($d=4$) are proved
uniform in the report and checked by \texttt{verify.py}. Weaker variants that
constrain only some shells are implied.
\item \textbf{The first conjunct.} The ``gap constant of the radius spectrum''
is not defined in the conjecture. We do not address that conjunct, and the
disproof does not need it.
\item \textbf{Linear.} Our codes are $\F$-subspaces. Linearity is part of the
Lean definition.
\end{itemize}

\section{Lean formalization}
The project \texttt{lean4/} (Lean 4.19.0, core library only) defines words as
\texttt{List Bool}, with weight \texttt{wt}, coordinatewise sum \texttt{xorW},
distance \texttt{dist} and the syndrome \texttt{synd i x}, the XOR of the labels
$i,i+1,\dots$ of the $1$-positions. A code is a predicate
\texttt{C : List Bool -> Prop}:
\begin{itemize}
\item \texttt{IsLinear n C}: all codewords have length $n$, $0\in C$, and $C$
is closed under \texttt{xorW}.
\item \texttt{MinDist C d}: all distinct codewords are at distance $\ge d$, and
some pair is at distance exactly $d$.
\item \texttt{CovRad n C r}: every word of length $n$ is within $r$ of $C$, and
some word is at distance $\ge r$ from all of $C$.
\item \texttt{IsQP n C}: \texttt{IsLinear n C} and, for some \texttt{d}, \texttt{MinDist C d} and \texttt{CovRad n C ((d-1)/2+1)}.
\end{itemize}
Lemmas \texttt{minDist\_unique} and \texttt{covRad\_unique} show these are
genuine (unique) values. The main theorems are:
\begin{itemize}
\item \texttt{SH\_quasiPerfect}: Theorem~\ref{thm:sh} for all $n$. It uses
\texttt{xor\_top\_lt} (Lemma~\ref{lem:bit}, proved via \texttt{Nat.testBit}).
\item \texttt{EH\_quasiPerfect}: the quasi-perfect part of
Theorem~\ref{thm:eh} for all $m\ge2$.
\item \texttt{EW\_quasiPerfect}: Proposition~\ref{prop:ew} without the shell
count. \texttt{no\_QP\_short}: Proposition~\ref{prop:short}.
\item \texttt{hamming\_not\_QP}: the Hamming codes of length $2^m-1$, $m\ge3$,
are not quasi-perfect. This shows that \texttt{IsQP} is not vacuous in the
other direction.
\item \texttt{QP\_spectrum}: \texttt{QPLen n} $\iff$ \texttt{2 <= n}.
\item \texttt{conjecture\_00000007964\_false}: $\neg$\,\texttt{LevenshteinClause}, where
\texttt{LevenshteinClause} is
$\exists L$ \texttt{: List Nat}, $\forall n$, \texttt{QPLen n} $\iff$ ($(\exists k,\ n+1=2^k)\lor n\in L$).
\item \texttt{onlyIf3\_false}, \texttt{onlyIf4\_false},
\texttt{onlyIf\_any\_false}, \texttt{onlyIfWeak\_false} and
\texttt{levenshteinAtLeast\_false}: the variants of Corollary~\ref{cor:main}.
\item Uniform-shell reading. The definitions are \texttt{DistTo C x r}
($r$ is the distance from $x$ to $C$) and \texttt{ShellUniform n C}: for any
two words of length $n$ at the same distance $r$ from $C$ and any $j$, there
are maps $f,g$ between the two $j$-shells with $g\circ f=\mathrm{id}$ and
$f\circ g=\mathrm{id}$ on them. Further, \texttt{UQPLen n} means that some
code satisfies both \texttt{IsQP n C} and \texttt{ShellUniform n C}. The
theorems are:
\begin{itemize}
\item \texttt{EW\_shellUniform}, for all $n$;
\item \texttt{SH5\_not\_shellUniform}, which shows the condition is a genuine
restriction;
\item \texttt{UQP\_spectrum}: \texttt{UQPLen n} $\iff$ \texttt{2 <= n};
\item \texttt{onlyIfUniform\_false}, and
\texttt{conjecture\_00000007964\_false\_uniform}:
$\neg$\,\texttt{LevenshteinClauseUniform}. This is the iff clause with
\texttt{UQPLen}.
\end{itemize}
\end{itemize}
The project contains no \texttt{sorry}, \texttt{native\_decide} or added
axioms. \texttt{\#print axioms} reports only \texttt{propext},
\texttt{Classical.choice} and \texttt{Quot.sound}. Not in Lean: the
shell-uniformity of $\mathrm{EH}(m)$ and $\mathrm{SH}(2^m-2)$
(Theorem~\ref{thm:uniform}(2),(3)), the shell counts of Theorem~\ref{thm:eh},
the repetition codes, and the dimension formula.

\section{Independent computer check}
\texttt{verify.py} (Python standard library) works by brute force. It lists
codewords explicitly, computes $d$ as the least nonzero weight, and computes
$\rho$ by a multi-source breadth-first search on the hypercube. It confirms:
\begin{itemize}
\item $\mathrm{SH}(n)$ for $3\le n\le 16$: $d=3$, with $\rho=2$ except at
$n\in\{3,7,15\}$, where $\rho=1$ (Hamming codes).
\item $\mathrm{EH}(m)$ for $m=2,3,4$: $d=4$, $\rho=2$, and the shell counts
of Theorem~\ref{thm:eh}.
\item $\mathrm{EW}(n)$ for $2\le n\le 14$, and the even repetition codes up to
length $16$.
\item An exhaustive enumeration of all binary linear codes of length $n\le 7$.
There are $1, 2, 5, 16, 67, 374, 2825, 29212$ of them, including $\{0\}$.
Quasi-perfect codes exist exactly for $n\ge 2$ in this range. Those with
$d\ge3$ exist for $n=4,5,6,7$ (5, 15, 107 and 1260 codes).
\item Shell uniformity, via coset weight distributions grouped by minimum
weight: $\mathrm{SH}(6)$, $\mathrm{SH}(14)$, $\mathrm{EH}(2)$,
$\mathrm{EH}(3)$, $\mathrm{EH}(4)$ and $\mathrm{EW}(n)$ for $n\le 12$ are
shell-uniform. Among $\mathrm{SH}(n)$ with $3\le n\le 14$, exactly
$n\in\{3,6,7,14\}$ are shell-uniform. $\mathrm{SH}(5)$ is not.
\item Lemma~\ref{lem:bit} for $k\le 12$.
\end{itemize}

\section{Reproduction}
\begin{verbatim}
cd lean4 && lake build      # Lean 4.19.0, no Mathlib
cd .. && python3 verify.py  # about one minute
pdflatex report.tex && pdflatex report.tex
\end{verbatim}

\end{document}
